\documentclass[11pt]{article}

\usepackage[margin=0.8in]{geometry}

% Typography
\usepackage{microtype}
\usepackage{newtxtext,newtxmath}

% Spacing + structure
\usepackage{titlesec}
\usepackage{parskip}
\usepackage{enumitem}
\usepackage[numbers,sort&compress]{natbib}
\usepackage{setspace}

% Tables
\usepackage{tabularx}
\usepackage{array}
\usepackage{ragged2e}

\pagestyle{empty}
\setlength{\parindent}{0pt}
\setlength{\parskip}{5pt}
\setstretch{1.06}

% Section spacing
\titlespacing*{\section}{0pt}{10pt}{4pt}
\titleformat{\section}{\normalsize\bfseries}{\thesection.}{0.6em}{}

% Lists
\setlist[itemize]{leftmargin=1.3em, topsep=2pt, itemsep=1pt, parsep=0pt}

% Timeline columns
\newcolumntype{L}[1]{>{\RaggedRight\arraybackslash}p{#1}}
\newcolumntype{Y}{>{\RaggedRight\arraybackslash}X}
\renewcommand{\arraystretch}{1.2}

\begin{document}

\begin{center}
{\Large\bfseries From Data Discovery to Analytics-Ready Data:\\
A Survey of Search, Preparation, and Integration in Data Lakes}\\[7pt]
{\normalsize Sumaiya Azad \quad $\cdot$ \quad u1494915 \quad $\cdot$ \quad University of Utah}
\end{center}

\vspace{2pt}

\section{Description}
Data lake pipelines are commonly described as three phases:
\emph{discovery and search} (finding relevant tables),
\emph{preparation and wrangling} (cleaning and structuring data for analysis),
and \emph{integration and matching} (aligning schemas and reconciling records).
This survey will review recent work across all three phases and will organize the literature by the
\emph{artifacts each phase produces} and the \emph{assumptions those artifacts encode}.

Across the literature, evaluation is often reported using \emph{phase-local} metrics that are not directly comparable across phases.
A key goal of this survey will be to summarize how papers define phase outputs, what datasets and benchmarks they use,
and which evaluation setups reflect utility for analytics-ready data.

\section{Method}
I will conduct a structured literature review consisting of:
\begin{itemize}
  \item a reproducible paper collection protocol (venues/years, search strings, inclusion/exclusion criteria);
  \item a taxonomy that maps papers to phase outputs and problem settings; and
  \item a synthesis comparing problem assumptions, produced artifacts, and evaluation practices across phases.
\end{itemize}

Anchor references will include 
LakeBench~\cite{10.14778/3659437.3659448} (benchmarks for discovery),
ALITE~\cite{10.14778/3574245.3574274} (table integration semantics),
and HAIPipe~\cite{10.1145/3588945} (human--machine preparation pipelines).

\section{Plans and Timeline}
\begin{tabularx}{\linewidth}{@{}L{2.6cm} @{\hspace{10pt}} Y@{}}
\textbf{Weeks 1--2} &
I will finalize the scope (years, venues, and tasks). I will write the review protocol: search strings and inclusion/exclusion criteria.
I will build an initial reading list from benchmark and system anchor works. \\[5pt]

\textbf{Weeks 3--4} &
I will screen and expand the corpus (including backward/forward citation chasing). I will assign papers to a first-pass taxonomy
for discovery and search, preparation and wrangling, and integration and matching. \\[5pt]

\textbf{Weeks 5--6} &
I will extract and normalize evaluation details (datasets, labels, metrics, and reported costs) using a consistent coding template. \\[5pt]

\textbf{Weeks 7--8} &
I will draft the survey synthesis by phase, emphasizing phase outputs and how papers define inputs/outputs, assumptions, and evaluation. \\[5pt]

\textbf{Weeks 9--10} &
I will finalize the write-up: taxonomy figure(s), comparison table(s) for datasets, and a concise set of open problems. \\
\end{tabularx}

\vspace{4pt}
{\small\bibliographystyle{abbrv}
\bibliography{ref}}

\end{document}
